% Abstract — drafted in Task 4
\begin{abstract}
Named Data Networking (NDN) replaces host-centric communication with content-centric named data retrieval, but its in-network caching and stateful forwarding expose routers to Interest Flooding Attacks (IFA) and Cache Pollution (CP) that evade conventional threshold-based detection. This paper presents a single unsupervised Isolation Forest detector, trained only on clean traffic, evaluated on two platforms across the same three topologies (tree, dumbbell, and DFN mesh): miniNDN, an emulation platform running the real NDN Forwarding Daemon, and ndnSIM, an ns-3 simulator. On miniNDN, the detector identifies the IFA attacker node at 91.2\,\%, 98.7\,\%, and 96.8\,\% across the three topologies, at a false-positive rate below 4.3\,\%, and confirms that PIT-based features are non-discriminating, correcting the textbook expectation of PIT exhaustion. Because CP leaves no usable signature in the emulation feature set, an attacker-rate sweep of 297 runs on ndnSIM locates each attack's detection floor, the minimum attacker rate at which reliable detection is sustained at a realistic 1--2\,\% false-positive rate. The sweep's central result: the two floors are an order of magnitude apart and topology-invariant, with CP, a volumetric attack, detected down to approximately 44\,Hz, and IFA, whose signature is qualitative rather than volumetric, detected down to approximately 1.8\,Hz. These findings indicate that the security-relevant limit on unsupervised NDN anomaly detection is not raw model accuracy but a mechanism-dependent detection floor set by how volumetric or qualitative each attack's traffic signature is.
\end{abstract}

\noindent\textbf{Keywords:} Named Data Networking; Interest Flooding Attack; Cache Pollution; Isolation Forest; anomaly detection; detection floor
